\section{Implementazione}

\subsection{PHP}

L'applicazione è sviluppata in PHP adottando una separazione netta tra presentazione e comportamento. Gli \emph{entry point} risiedono in \texttt{public/} (ad es.\ \texttt{index.php}, \texttt{prodotti.php}, \texttt{prodotto.php}, \texttt{login.php}, \texttt{signup.php}, \texttt{logout.php}, \texttt{area\_personale.php}) e hanno il solo compito di caricare l'autoloader e di delegare al servizio pertinente la costruzione della pagina. La logica applicativa è organizzata in \texttt{src/}.

Il rendering è orchestrato dalla classe \texttt{App\textbackslash Core\textbackslash PageBuilder}, che inizializza la sessione, verifica la disponibilità dei servizi essenziali (autenticazione e database) e compone l'output finale integrando \emph{head}, \emph{header}, \emph{footer} e contenuto specifico della pagina. Le parti comuni dell'interfaccia sono costruite dai \emph{builder} in \texttt{App\textbackslash View} (\texttt{HeadBuilder}, \texttt{HeaderBuilder}, \texttt{FooterBuilder}), richiamati dal \texttt{PageBuilder} o dai servizi quando necessario. Questa scelta elimina duplicazioni tra pagine e consente modifiche centralizzate del layout.
 

La presentazione è definita da template HTML collocati in \texttt{src/html/}, privi di logica PHP. La sostituzione dei segnaposto è affidata al micro--motore \texttt{App\textbackslash Core\textbackslash Template}, che gestisce placeholder del tipo \texttt{\{\{ id \}\}} e blocchi \texttt{\{\{ id \}\} \dots \{\{/id\}\}}, inclusi casi annidati e strutture array/oggetto. 

L'accesso ai dati è incapsulato in \texttt{App\textbackslash Core\textbackslash Database}, che fornisce una connessione \texttt{mysqli} gestita come singleton. Le interrogazioni al database utilizzano \emph{prepared statements} per prevenire SQL injection; per disaccoppiare persistenza e presentazione impieghiamo oggetti di trasferimento dati (DTO), come \texttt{App\textbackslash Core\textbackslash Model\textbackslash ProductDTO} e \texttt{UserDTO}. Quando è necessario costruire elenchi filtrabili e paginati, le classi di supporto in \texttt{App\textbackslash Core\textbackslash Filter} e \texttt{App\textbackslash Core\textbackslash Pagination} compongono in maniera sistematica condizioni, limiti e offset, mantenendo le regole di business fuori dalle query ad hoc.

L'autenticazione è gestita da \texttt{App\textbackslash Service\textbackslash AuthService}, che utilizza \texttt{password\_hash()} e \\\texttt{password\_verify()} per il trattamento delle credenziali e si occupa della persistenza di sessione. I flussi di login e registrazione, coordinati rispettivamente da \texttt{LoginService} e \texttt{SignupService}, validano gli input lato server, impostano messaggi di errore e forniscono alla vista un modello coerente con lo stato dell'operazione. Nei form sensibili, in particolare nell'area riservata governata da \texttt{AreaPersonaleService}, introduciamo token anti--CSRF generati e verificati sul server; i messaggi di validazione sono esposti tramite elementi con \texttt{role="alert"} e \texttt{aria-live="assertive"} per una corretta fruizione con tecnologie assistive.

\subsection{Javascript}

Adottiamo un approccio di \emph{progressive enhancement}: la pagina e i form funzionano anche senza script (invio tradizionale), mentre con JavaScript abilitiamo l’aggiornamento in-place della lista prodotti in risposta ai filtri. Al cambio dei controlli serializziamo i parametri con \texttt{FormData}/\texttt{URLSearchParams} e inviamo una richiesta \texttt{fetch} a \texttt{/prodotti.php}, marcando \texttt{\#items} con \texttt{aria-busy} e usando l’header \texttt{X-Requested-With: fetch} per consentire al server di restituire, quando possibile, solo il frammento HTML della lista; la risposta, sia frammento sia pagina intera, viene parsata e reinserita in \texttt{\#items}. 

Le \emph{validazioni} sono progettate in modo \emph{server-authoritative}: la verifica effettiva dei parametri avviene sempre lato server e l’esito non dipende da ciò che succede nel browser. Quando un input non supera i controlli, il server restituisce un frammento HTML accessibile con i messaggi di errore, esposto con \texttt{role="alert"} e \texttt{aria-live="assertive"}, così che l’utente sia informato in modo chiaro e compatibile con le tecnologie assistive. In parallelo, gli \emph{script di validazione lato client} affiancano il flusso senza sostituirlo: si agganciano agli eventi \texttt{input}/\texttt{change}/\texttt{submit}, normalizzano i valori, marcano i campi con \texttt{aria-invalid} e collegano i messaggi con \texttt{aria-describedby}. In questo modo l’esperienza resta fluida (aggiornamenti in-place quando possibile).